\section{Decay survives denominator clearing}\label{sec:decay}

The prime number theorem gives $d_n^{1/n}\to\ep$, hence also
$\delta_n^{1/n}\to\ep$. The multiplier $\delta_n^{365}$ therefore grows
at exponential rate $\ep^{365}$. The following bound shows that the
forms decay fast enough to compensate.

\begin{proposition}\label{prop:decay}
For $b=1,2,3$, along positive even $n$,
\[
 \limsup_{n\to\infty}|\delta_n^{365}r_{n,b}|^{1/n}<\frac34.
\]
\end{proposition}

\begin{proof}
The positive tail begins just beyond the last zero $7n$. For fixed
$x>0$, its logarithmic growth rate at distance $nx$ from that zero is
\[
 \Phi(x):=\lim_{\substack{n\to\infty\\n\text{ even}}}
                 \frac1n\log R_n\bigl(n(7+x)\bigr).
\]
Thus, for fixed $x>0$, $R_n(7n+nx)=\exp(n\Phi(x)+o(n))$.
We first find the maximum of $\Phi$, then use uniform estimates to
bound the whole tail.

The product identity
$\prod_{j=0}^{N-1}(z+j)=\Gamma(z+N)/\Gamma(z)$ gives
\begin{align*}
 P_n(7n+t)
   &=t\prod_{\theta\in\Theta}
          \frac{\Gamma(t+18n+\theta)}{\Gamma(t+\theta)},\\
 \prod_{h=0}^{n}(7n+t+4h)
   &=4^{n+1}\frac{\Gamma((11n+t)/4+1)}{\Gamma((7n+t)/4)}.
\end{align*}
Substitute these expressions and $C_n$ from \eqref{eq:rn} into
$\log R_n(7n+t)$, set $t=nx$, and apply Stirling's formula. The limit
exists and equals
\begin{align}
 \Phi(x)={}&36\log2+54\log3
     +4\bigl((x+18)\log(x+18)-x\log x\bigr)\notag\\
 &-\frac{365}{4}\bigl((x+11)\log(x+11)-(x+7)\log(x+7)\bigr)
     +365\log4,\label{eq:profile}
\end{align}
with continuous extension to $x=0$ obtained by setting $x\log x=0$.
The $n\log n$ terms and the linear terms from Stirling's formula cancel.
The factor $4^{365n}$ in $C_n$ contributes the final term $365\log4$.

To locate the maximum, put
\begin{equation}\label{eq:f}
 4\Phi'(x)=\log f(x),\qquad
 f(x)=\left(\frac{x+18}{x}\right)^{16}
      \left(\frac{x+7}{x+11}\right)^{365}.
\end{equation}
Its logarithmic derivative is
\[
 (\log f)'(x)=-\frac{288}{x(x+18)}
                  +\frac{1460}{(x+7)(x+11)}.
\]
The numerator over a common positive denominator is
$1172x^2+21096x-22176$, which has one positive root. Thus $\log f$
first decreases and then increases. It tends to $+\infty$ as $x\to0$
and to zero from below as $x\to\infty$, with leading term $-1172/x$.
Consequently $f(x)=1$ has exactly one positive solution $x_0$, where
$\Phi$ attains its global maximum. Using $\Phi'(x_0)=0$ to simplify
\eqref{eq:profile} gives
\begin{equation}\label{eq:G}
 G:=\ep^{\Phi(x_0)}
  =2^{36}3^{54}(x_0+18)^{72}4^{365}
      \left(\frac{(x_0+7)^7}{(x_0+11)^{11}}\right)^{365/4}.
\end{equation}

We now bound the sum in \eqref{eq:positive-bound}. For fixed $A$, the
gamma-product formula gives
\[
 \log R_n(7n+t)=n\Phi(t/n)+O_A(\log n)
                  \quad(1/3\leq t\leq An),
\]
uniformly on that interval. The estimate remains valid for small $t$ because
$\log\Gamma(y)=y\log y-y+O(\log(y+1)+1)$ for $y$ bounded below by
a fixed positive constant. Each grid has $O_A(n)$ terms in this interval,
so the sum over the interval has $n$th-root limsup at most $G$.

For $t\geq An$ with $A\geq18$, every numerator factor is at most $2t$
and every denominator factor is at least $t$. Therefore
\[
 R_n(7n+t)\leq C_n2^{72n+1}t^{-M_n},
 \qquad M_n=293n+364.
\]
The grid has spacing $1/b$. For the sum of $t^{-M_n}$ over this grid,
we bound the first term by $(An)^{-M_n}$ and the rest by an integral:
\begin{align*}
 \sum_{j/b\geq(7+A)n}R_n(j/b)
 &\leq C_n2^{72n+1}
       \left((An)^{-M_n}+b\int_{An}^{\infty}t^{-M_n}\,dt\right)\\
 &=C_n2^{72n+1}(An)^{-M_n}
       \left(1+\frac{bAn}{M_n-1}\right).
\end{align*}
Stirling's formula gives
\[
 C_n=4^{365n}K^n(n+1)^{365}(n!)^{293}
     =\bigl(4^{365}K\ep^{-293}\bigr)^n
        n^{293n}\exp\bigl(O(\log n)\bigr).
\]
The factor $n^{293n}$ cancels the same power in $(An)^{-M_n}$.
The remaining powers of $n$ do not affect the $n$th-root limit.
Taking $n$th roots in the tail bound therefore gives a limsup of at most
\[
 4^{365}K2^{72}\ep^{-293}A^{-293}.
\]
This tends to zero as $A\to\infty$. We conclude that
\begin{equation}\label{eq:unscaled-decay}
 \limsup_{\substack{n\to\infty\\n\text{ even}}}|r_{n,b}|^{1/n}\leq G.
\end{equation}

To prove $\ep^{365}G<3/4$, set $u=3/5000$.
Exact rational arithmetic gives
\begin{equation}\label{eq:fu}
 \left(\frac{u+18}{u}\right)^{16}
 \left(\frac{u+7}{u+11}\right)^{365}<1,
\end{equation}
so $x_0<u$. Also $(x+7)^7/(x+11)^{11}$ is decreasing for $x\geq0$,
since its logarithmic derivative is $-4x/((x+7)(x+11))$. Finally,
the positive exponential series gives
\[
 \ep<\sum_{k=0}^5\frac1{k!}
              +\frac{1/6!}{1-1/7}
       =\frac{163}{60}+\frac7{4320}<\frac{68}{25}.
\]
Consequently
\begin{align}
 (\ep^{365}G)^4
 &<\left(\frac{68}{25}\right)^{1460}2^{144}3^{216}
       \left(18+\frac3{5000}\right)^{288}4^{1460}
       \left(\frac{7^7}{11^{11}}\right)^{365}\notag\\
 &<\frac{81}{256}=\left(\frac34\right)^4.
 \label{eq:rate-certificate}
\end{align}
The rational expression in the first line is approximately $0.2761910348$;
its exact comparison with $81/256$ is verified below. Combining
\eqref{eq:unscaled-decay} with $\delta_n^{1/n}\to\ep$ proves the proposition.
\end{proof}

The comparisons \eqref{eq:fu} and \eqref{eq:rate-certificate}, including
the bound for $\ep$, can be reproduced with the following exact calculation:
\begin{quote}\small
\begin{verbatim}
from fractions import Fraction as F
u = F(3, 5000)
assert ((u+18)/u)**16 * ((u+7)/(u+11))**365 < 1
assert F(163,60) + F(7,4320) < F(68,25)
U = (F(68,25)**1460 * 2**144 * 3**216 * (18+u)**288
     * 4**1460 * (F(7)**7/F(11)**11)**365)
assert U < F(81,256)
\end{verbatim}
\end{quote}
